\section{Theoretical Framework: Error Accumulation}

\subsection{Markov Chain Model}

We model reasoning as a Markov process where each token generation step can introduce errors.

\textbf{State space:} $S = \{\text{correct}, \text{wrong}\}$

\textbf{Transition probabilities:}
\begin{align}
P(\text{correct} \to \text{correct}) &= 1 - \epsilon(t) \\
P(\text{correct} \to \text{wrong}) &= \epsilon(t) \\
P(\text{wrong} \to \text{correct}) &= \delta(t) \\
P(\text{wrong} \to \text{wrong}) &= 1 - \delta(t)
\end{align}

where $\epsilon(t)$ is the error rate at step $t$, and $\delta(t)$ is the self-correction rate.

\subsection{When Does Overthinking Occur?}

\begin{theorem}[Overthinking Condition]
If $\epsilon(t) > \delta(t)$ for all $t > t_0$, then accuracy monotonically decreases after step $t_0$.
\end{theorem}

\begin{proof}
Let $p_t$ be the probability of being in the correct state at step $t$.
Then:
$$p_{t+1} = p_t(1-\epsilon(t)) + (1-p_t)\delta(t)$$

For overthinking to occur, we need $p_{t+1} < p_t$:
\begin{align}
p_t(1-\epsilon(t)) + (1-p_t)\delta(t) &< p_t \\
\delta(t) - p_t\delta(t) &< p_t\epsilon(t) \\
\delta(t) &< p_t(\epsilon(t) + \delta(t))
\end{align}

When $p_t \approx 1$ (already correct), this simplifies to $\delta(t) < \epsilon(t)$.
\end{proof}

\textbf{Interpretation:} Overthinking occurs when the error introduction rate exceeds the self-correction rate.
